\documentclass[a4paper,12pt]{letter}
\usepackage{geometry}
\usepackage{graphicx}
\geometry{a4paper, margin=0.8in}

\begin{document}

% Letterhead with logo aligned at top (saves vertical space)
\noindent
\begin{minipage}[t]{0.65\textwidth}
University of Graz \\
Wegener Centre for Climate and Global Change \\[0.5em]
Graz, Austria, \today
\end{minipage}%
\hfill
\begin{minipage}[t]{0.3\textwidth}
\raggedleft
\includegraphics[width=\textwidth]{Figures_supp/Logo_Universitaet-Graz_Schriftzug_1c.jpg}
\end{minipage}

\vspace{1em}

\noindent
To: The Editors, Nature Communications Physics

\vspace{0.5em}

\noindent
\textbf{Re: Submission of Manuscript titled "Guided rewiring of social networks reduces polarization and accelerates collective action"}

\vspace{0.5em}

\noindent
Dear Editors,

\vspace{0.5em}

On behalf of my co-authors, I am pleased to submit our article "Guided rewiring of social networks reduces polarization and accelerates collective action" for your consideration.

Evidence shows that contemporary algorithmic design, as utilized by very large online platforms (VLOPS) such as X (formerly Twitter), and Facebook may polarize participants in digital common spaces and hinder cooperative consensus formation. Previous work has investigated the effect of commercially deployed link-recommendation algorithms on social network inequalities and, separately, the effect of idealized link-recommendations on opinion dynamics. A gap exists, however, in the combination of these commercially deployed algorithms with classical opinion dynamics models.

We formulate two rewiring algorithms to reduce polarization and compare these against two commonly used rewiring algorithms employed by VLOPS across real and idealized network topologies. We show that one of the commercial algorithms, Who To Follow, increased polarization across empirical and synthetic networks and prohibited cooperative consensus formation.

Our work contributes to the literature in that it (1) combines an opinion dynamics model with commercial-use topological rewiring algorithms derived from Who To Follow and node2vec, (2) tests this on real-world network structures, (3) provides algorithms that per our results are better than these at reducing polarization, and (4) includes directed and undirected networks. Our results inform strategies which improve the capacity for online spaces to accelerate collective action on urgent global issues.

Our work will be of interest to readers of Nature Communications Physics as it addresses how to resolve the polarization that currently impedes effective collective action on large-scale challenges in a methodological space at the intersection of network theory, complex systems, statistical physics, and computational social science.

Please do not hesitate to contact me if you have any questions or if you need any further information.

\vspace{1em}

\noindent
Yours sincerely,

\vspace{1em}

\noindent
Jordan Everall, MSc \\
PhD candidate \\
Wegener Centre for Climate and Global Change \\
University of Graz \\
Email: jordan.everall@uni-graz.at

\end{document}
